\section{Evaluation}
\label{sec:evaluation}

\subsection{Research Questions}

The evaluation is organized around three research questions. Source
characterization selects workloads and source oracles, and same-oracle
comparisons plus boundary evidence answer the RQs.
The unit of evidence is a complete source-oracle matrix, not a catalog of
one-off baselines or mechanism smoke tests.

\begin{description}
\item[RQ1] Can a narrow VFS name-resolution extension express real
state-dependent path-view policies without taking over filesystem semantics?
\item[RQ2] What is the cost of putting programmable policy on the VFS
name-resolution path compared with a feature-equivalent FUSE policy
implementation?
\item[RQ3] Does \namei provide a narrower verifier-bounded, fail-closed
ownership boundary than building a custom or stackable filesystem when the
needed policy is only name resolution?
\end{description}

\subsection{Experimental Setup}

All \namei mechanism runs execute in KVM~\cite{linux_kvm} with the modified
kernel and the real \code{cgroup/namei_ext} attach path. Make targets own
workload execution, comparison construction, result collection, and report
generation. Each run records the machine, kernel, image, policy, workload
inputs, stdout/stderr, dmesg, per-operation lookup/readdir events, and
comparison-specific setup observations.

Correctness gates every row. A source oracle is the correctness condition
extracted from the original source system. A same-oracle mechanism comparison
requires the executed \namei and feature-equivalent FUSE rows to pass that same
condition, not weaker mechanism-specific tests. RQ3 boundary accounts use the
same source oracle to compare responsibilities and containment rather than to
claim an executed custom/stackable correctness row. Reproduced or replayed
source oracles can become evaluation rows. Source inspection and paper-derived
evidence remain context unless they are backed by a real KVM oracle.

\begin{table}[t]
\centering
\scriptsize
\begin{tabular}{@{}L{0.18\linewidth}L{0.26\linewidth}L{0.2\linewidth}L{0.24\linewidth}@{}}
\hline
Experiment & Source oracle & \namei row & Required comparisons \\
\hline
Agent workspace lifecycle & complete AgentFS-derived branch/checkpoint, COW,
whiteout, symlink, and final-tree behavior & full KVM
\code{cgroup/namei_ext} run over the complete fixed workspace lifecycle,
attributing only lookup/readdir decisions to \namei &
feature-equivalent FUSE policy for RQ2, plus source-system and custom/stackable
ownership boundary account for RQ3 \\
Environment/cache transition & MEnvData-SWE and SWE-Factory-Gym build/test
rows with hit, miss, stale, corrupt, and epoch-update states & full KVM
\code{cgroup/namei_ext} run over the fixed cache-state machine, with
file-object evidence requiring final-file target selection &
feature-equivalent FUSE cache-policy view for RQ2, plus source-native and
custom/stackable cache ownership boundary account for RQ3 \\
Service/config transition & concrete source oracle only when lookup-time
object selection changes service-visible behavior & evaluated only if a source
oracle requires lookup-time object selection & FUSE and boundary comparisons
only for an included source oracle \\
\hline
\end{tabular}
\caption{Complete experiment shape. The evaluation has two primary matrices.
Service/config is evaluated only if a source oracle requires
lookup-time object selection.}
\label{tab:experiment-shape}
\end{table}

\paragraph{Measurement protocol.}
For each final reported run, the artifact records hardware and kernel
configuration, kernel commit, KVM image digest, filesystem configuration, run
count, warmup rule, confidence-interval method, cache-hot/cache-cold protocol,
feature-equivalent FUSE implementation, and custom or stackable filesystem
boundary account for each reportable workload.

\subsection{RQ1: Expressiveness And Sufficiency}

RQ1 asks whether the VFS boundary can express path-view transitions extracted
from source oracles while preserving lower-filesystem semantics. A reportable
workload must run the complete source lifecycle or fixed replay needed by its
oracle, while \namei evidence is limited to lookup and readdir effects.
Runtime orchestration, storage state, synthetic contents, write handling, and
source-system lifecycle behavior remain delegated to the source system or
lower filesystem.

\begin{table}[t]
\centering
\scriptsize
\begin{tabular}{@{}L{0.2\linewidth}L{0.24\linewidth}L{0.22\linewidth}L{0.24\linewidth}@{}}
\hline
Source/state & Path operation/effect & Oracle & Delegated behavior/boundary \\
\hline
AgentFS-derived workspace branch/checkpoint~\cite{agentfs_repo} & lookup/readdir select or hide branch-visible lower objects during the lifecycle & final tree, lookup/readdir coherence, lower-filesystem permissions/writes & complete lifecycle run with COW, checkpoint, sandbox orchestration, and writes kept with the source runtime or lower filesystem \\
Environment/cache epoch~\cite{swe_factory_repo,menvagent_repo,swe_rebench_v2_repo} & lookup selects verified local or canonical fallback, or hides stale/corrupt local objects as reported by source cache metadata/validation & source build/test result plus stale/corrupt non-use & representative workload with cache population, content validation, and build/test execution kept with source/lower filesystem, with final-file target support required for file-object selection \\
Service/config epoch~\cite{kubernetes_projected_volumes,kubernetes_configmaps,kubernetes_secrets} & lookup selects epoch object at stable service path & service/source oracle for lookup-time object choice & included when a concrete oracle is chosen, with reload protocol and data path outside the hook \\
\hline
\end{tabular}
\caption{Workload evidence contract. Only the path operation/effect column is
\namei evidence, and delegated behavior remains outside the extension
boundary.}
\label{tab:workload-evidence}
\end{table}

Each representative workload must pass its source oracle in KVM before
inclusion, preserve lower-filesystem permission/write/data-path behavior,
exercise lookup/readdir coherence, and record an operation-weighted trace.

\subsection{RQ2: Cost Versus FUSE}

RQ2 compares \namei with a feature-equivalent FUSE policy service under the same
source oracle, isolating the request-path cost of moving the same decision out
of VFS name resolution. For each final reported workload, the FUSE row must
implement the same source policy and pass the same oracle before any overhead
is interpreted. Non-name-resolution lifecycle
behavior is delegated identically in both rows, so the comparison isolates the
placement of the lookup/readdir policy rather than giving either mechanism a
different runtime.

The reported metrics are pass-through and action-specific \namei overhead,
lookup/open/stat/access/exec/readdir latency, cache-hot and cache-cold behavior,
macro workload runtime, operation-weighted invocation rate, matched action mix,
and FUSE daemon/request-path cost. No-hook and pass-through controls separate
fixed hook cost from policy cost. The same-oracle FUSE rows show the cost of
placing equivalent policy in a filesystem service.

\subsection{RQ3: Safety And Boundary Versus Custom Or Stackable Filesystems}

RQ3 asks whether verifier-bounded policy plus kernel validation gives \namei a
narrower verifier-bounded, fail-closed ownership boundary than writing a custom
or stackable filesystem for policies that only need name resolution. Here,
policy-output safety and containment mean that the eBPF verifier bounds policy
execution, the kernel validates names, action types, target identifiers, and
selection scope, malformed decisions fail closed, and lower filesystems still
enforce permissions, data operations, writes, and persistence. The claim is not
that \namei proves full filesystem correctness or replaces mechanisms needed
for richer semantics. The comparison is not whether custom filesystems are
expressive. Instead, it asks which filesystem responsibilities the developer
must implement or operate for the same source oracle.

The boundary account records same-oracle KVM pass/fail, preserved
lower-filesystem permissions and writes, malformed or unsupported decision
containment through \code{cgroup/namei_ext}, filesystem methods and storage
semantics the developer must implement, privileged code surface, policy code
size, state responsibility, and whether the source policy needs behavior beyond
name resolution. Custom or stackable filesystems place policy in filesystem
methods, while \namei keeps policy in verified eBPF and leaves data, write, and
persistence semantics with the lower filesystem.

RQ3 supports a narrower verifier-bounded, fail-closed ownership boundary only
if the same oracle passes, invalid decisions fail visibly through the real
attachment path, and the comparison shows that a custom or stackable filesystem
would own filesystem methods, daemon state, storage semantics, or data-path
behavior that \namei leaves to the VFS and lower filesystem.

For the Agent workspace matrix, the boundary account compares the complete
source lifecycle's name-resolution decisions with AgentFS, BranchFS, YoloFS,
and custom or stackable filesystem ownership. For the environment/cache matrix,
it compares the cache-state name-resolution policy with the source-native
environment mechanism, feature-equivalent FUSE ownership, and any custom or
stackable filesystem responsibility needed to own cache metadata, stale/corrupt
state, or data-path behavior.

\subsection{Interpretation Scope}

Rows that require final-file target selection, including environment/cache
file-object evidence, are outside the evidence reported here until that support
exists. Service/config is included only when a concrete lookup-time source
oracle exists. RQ answers come from reviewed KVM source-oracle runs,
feature-equivalent FUSE comparisons, and ownership-boundary evidence.
Materialized namespace mechanisms, eBPF LSM hooks, and native production
mechanisms remain related work and background context unless a selected source
workload makes one of them part of the source behavior being evaluated.
